The restriction maps are identities when both opens contain $P$ and zero otherwise. On a cover, all sections over members containing $P$ must agree, since their overlaps contain $P$; this gives unique gluing. Thus $i_P(A)$ is a sheaf.

If $Q\in\overline{\{P\}}$, every neighborhood of $Q$ contains $P$, so the stalk is the direct limit of copies of $A$ with identity maps, hence $A$. Otherwise there is a neighborhood of $Q$ missing $P$, and the stalk is zero.

Every nonempty open subset of $\overline{\{P\}}$ contains $P$, so this closed subspace is irreducible, and its constant sheaf has value $A$ on every nonempty open subset. Also an open subset of $X$ meets $\overline{\{P\}}$ exactly when it contains $P$. Therefore its direct image under $i:\overline{\{P\}}\to X$ agrees with $i_P(A)$ on sections and restrictions.
